%%%%%%%%%%%%%%%%%%%%%%%%%%%%%%%%%%%%%%%%%%%%%%%%%%%%%%%%%%%%%%%%%%%%%%%%%%%%%%
%% macros.tex -- notation shortcuts and the open-slot ("FILL") machinery
%%%%%%%%%%%%%%%%%%%%%%%%%%%%%%%%%%%%%%%%%%%%%%%%%%%%%%%%%%%%%%%%%%%%%%%%%%%%%%

%%%%%%%%%%%%%%%%%%%%%%%%%%%%%%%%%%%%%%%%%%%%%%%%%%%%%%%%%%%% open slots %%%%%
%
% Every number that depends on the running campaign is marked with \FILL or
% \FILLBOX. Both are listed, with page numbers, by \listoffills at the end of
% the document.
%
% Set \showfillfalse (below) to make every marker vanish from the PDF --
% useful for checking page count against the venue limit, since the markers
% themselves consume space.

\newif\ifshowfill
\showfilltrue          % <-- flip to \showfillfalse for a clean, marker-free PDF

\definecolor{fillbg}{HTML}{FFF3CD}
\definecolor{fillrule}{HTML}{D39E00}
\definecolor{filltext}{HTML}{7A5C00}

% Record every slot into the .fil auxiliary file so \listoffills can print it.
\newcommand{\recordfill}[1]{%
  \addcontentsline{fil}{figure}{\protect\numberline{}\small #1}%
}

% --- inline slot: \FILL{what is missing} ------------------------------------
% The tag is boxed (short, never needs to break); the text is plain coloured
% prose so it wraps normally. Do not put the text inside \colorbox -- an
% unbreakable box produces enormous overfull lines.
\newcommand{\FILL}[1]{%
  \recordfill{#1}%
  \ifshowfill
    \colorbox{fillbg}{\textcolor{filltext}{\scriptsize\textsf{\bfseries FILL}}}%
    \,\textcolor{filltext}{\small #1}%
  \fi
}

% --- paragraph-level slot: \begin{FILLBOX} ... \end{FILLBOX} ----------------
\newenvironment{FILLBOX}[1][open slot]{%
  \recordfill{#1}%
  \ifshowfill
    \begin{tcolorbox}[
      breakable, enhanced, sharp corners,
      colback=fillbg, colframe=fillrule, boxrule=0.4pt,
      left=5pt, right=5pt, top=4pt, bottom=4pt,
      before skip=6pt, after skip=6pt,
      fontupper=\small\color{filltext}]%
    \textsf{\bfseries [FILL]}\ \ignorespaces
  \else
    \setbox0\vbox\bgroup
  \fi
}{%
  \ifshowfill
    \end{tcolorbox}%
  \else
    \egroup
  \fi
}

% --- the index of open slots -------------------------------------------------
\makeatletter
\newcommand{\listoffills}{%
  \section*{Open slots in this draft}%
  \small
  \@starttoc{fil}%
}
\makeatother

% Placeholder for a number not yet measured, inside tables and running text.
\newcommand{\tbd}{\textcolor{fillrule}{\textbf{--}}}

%%%%%%%%%%%%%%%%%%%%%%%%%%%%%%%%%%%%%%%%%%%%%%%%%%%%%%%%%%%%%%% notation %%%%

\newcommand{\psifield}{\psi}
\newcommand{\Rgrid}{R}
\newcommand{\Zgrid}{Z}
\newcommand{\relltwo}{\mathrm{relL2}}
\newcommand{\Ip}{I_{\mathrm{p}}}
\newcommand{\rzero}{r_{0}}
\newcommand{\gs}{Grad--Shafranov}
\newcommand{\GS}{Grad--Shafranov}

% levels and harnesses, so they are typeset consistently everywhere
\newcommand{\Lzero}{\textsf{L0}}
\newcommand{\Lone}{\textsf{L1}}
\newcommand{\Ltwo}{\textsf{L2}}
\newcommand{\Lthree}{\textsf{L3}}
\newcommand{\hn}[1]{\textsf{#1}}

% code-ish inline
\newcommand{\code}[1]{\texttt{\small #1}}

\newcommand{\eg}{e.g.\xspace}
\newcommand{\ie}{i.e.\xspace}
